\documentclass[12pt]{article}

\usepackage{amsmath}
\usepackage{amssymb}
\usepackage{graphicx}
\graphicspath{{figures/}}
\usepackage{hyperref}
\usepackage[margin=1in]{geometry}
\usepackage{tcolorbox}
\tcbuselibrary{breakable,skins}
\usepackage{natbib}
\bibliographystyle{aasjournal}

\title{Restriction of a Density Field and the 1LPT Displacement}
\author{}
\date{\today}

\begin{document}

\maketitle

\begin{tcolorbox}[colback=yellow!10!white, colframe=orange!75!black,
                  title={\bfseries Disclaimer}]
These notes were compiled by an AI assistant (Claude, Anthropic) and have
\textbf{not been reviewed by a human domain expert}.  Cross-check key
equations and conclusions before relying on them.
\end{tcolorbox}

\section{Setup}

Take a density field $\delta_{2N}$ sampled on a periodic cubic grid with
$2N$ points per side and spacing $\Delta x_f = L/(2N)$. The
\emph{restriction} operator $R$ replaces each $2\times2\times2$ block of
fine cells (or $2\times2$ in 2D) by its arithmetic mean, producing
$\delta_N = R\,\delta_{2N}$ on a grid of $N$ points per side and spacing
$\Delta x_c = L/N = 2\Delta x_f$. Each output cell is the volume average
of the eight (in 3D; four in 2D) fine cells inside it, which makes the
restriction \emph{locally} mass conserving: the mass in any coarse cell
$[I,J,K]$ equals the sum of the masses in its child fine cells,
\begin{equation}\label{eq:local-mass}
  (1 + \delta_{N,IJK})\,\Delta x_c^3
  \;=\;\!\!\!\sum_{(i,j,k)\in\,\mathrm{block}(I,J,K)}\!\!\!
        (1 + \delta_{2N,ijk})\,\Delta x_f^3,
\end{equation}
where the block contains the eight fine cells $\{(2I+a,2J+b,2K+c)\}_{a,b,c\in\{0,1\}}$.
Because $\Delta x_c^3 = 8\Delta x_f^3$, equation~(\ref{eq:local-mass})
reduces to $\delta_{N,IJK} = \tfrac{1}{8}\sum_{\rm block}\delta_{2N}$,
i.e.\ block averaging. Summing (\ref{eq:local-mass}) over all coarse
cells recovers global mass conservation as an immediate consequence.

The questions addressed here:
\begin{enumerate}
\item How is the Fourier transform of $\delta_N$ related to that of
      $\delta_{2N}$?
\item In first-order Lagrangian perturbation theory (1LPT), how does
      the particle displacement field $\boldsymbol{\Psi}^{(1)}$ change
      when computed from $\delta_N$ rather than $\delta_{2N}$?
\item How does the deformation tensor $\partial_i\Psi_j^{(1)}$ --- and
      the second-order LPT (2LPT) source built from it --- differ
      between the two grids?
\item What goes wrong if one instead restricts the displacement
      potential $\phi^{(1)}$ and then differentiates?
\end{enumerate}

All numerical illustrations are 2D for ease of visualization, with
$N=64$, $2N=128$, and box length $L=1$. The fine-grid density
$\delta_{2N}$ is a synthetic Gaussian random field generated by drawing
a unit-variance white noise field $\mu(\mathbf q)\sim\mathcal{N}(0,1)$
in real space, taking its Fourier transform, multiplying by $\sqrt{P(k)}$
with a toy power spectrum $P(k)\propto k^n$ (we use $n=-2$ in the
figures below, chosen as the proper 2D analog of the high-$k$ slope of
$\Lambda$CDM --- see the dimensional argument in
\S\ref{sec:cdm}), setting the
$\mathbf k=\mathbf 0$ mode (the spatial mean) to zero, and transforming
back. This is a placeholder spectrum chosen for illustration --- it is
not the physical $\Lambda$CDM matter power spectrum --- but it has
enough small-scale power for the window and alias effects of restriction to
be visible (their absolute size depends on the slope $n$ chosen --- a
shallower spectrum makes them more pronounced). Zeroing the $\mathbf k=\mathbf 0$ mode makes
$\delta_{2N}$ have exactly zero mean on the grid by construction.

\section{Restriction in real and Fourier space}\label{sec:fourier}

Restriction is a linear operation. To see this in Fourier space,
think of it in two pieces: a \emph{convolution} of the fine-grid
samples with a small averaging kernel, followed by a $2{:}1$
\emph{subsampling} that keeps every other point. (For a self-contained
review of the Fourier-analysis tools used in this section ---
sampling, the Nyquist limit, aliasing, the convolution theorem ---
see \texttt{notes/fft.tex}.)

The discrete convolution is the grid analogue of the familiar
continuous convolution
$(h*f)(x)=\int h(x-x')f(x')\,dx'$: replace the integral by a sum and
the continuous kernel $h(x)$ by a finite list of weights $h_n$. For
each output position $m$, the convolved value is a weighted sum of
nearby input samples,
\begin{equation}
  y_m \;=\; \sum_n h_n\,\delta_{2N,m+n}.
\end{equation}
The list $\{h_n\}$ is called a \emph{kernel}; for our restriction
operation it is the simplest possible smoothing kernel, an equal-weight
average over two adjacent fine cells.

In 1D, define the kernel\footnote{The integer $n$ here is just a
position label on the fine grid, in units of the fine spacing $\Delta x_f$:
$n=0$ refers to the sample at position $x=0$, $n=1$ to the sample at
$x=\Delta x_f$, and so on.}
\begin{equation}
  h_n \;=\;
  \begin{cases}
    \tfrac12 & n=0,\\
    \tfrac12 & n=1,\\
    0 & \text{otherwise}.
  \end{cases}
\end{equation}
Plugging this into the convolution formula gives
$y_m = \tfrac12\bigl(\delta_{2N,m} + \delta_{2N,m+1}\bigr)$,
the moving average of two adjacent fine samples (defined on the same
fine grid as $\delta_{2N}$). Subsampling means we then keep only the
even-indexed entries:
\begin{equation}\label{eq:discrete-restrict}
  \delta_{N,I} \;=\; y_{2I}
              \;=\; \tfrac12\bigl(\delta_{2N,2I} + \delta_{2N,2I+1}\bigr).
\end{equation}
Equation~(\ref{eq:discrete-restrict}) is exactly $2{:}1$ block averaging:
the coarse cell with index $I$ is the mean of the two fine cells with
indices $2I$ and $2I+1$. The factor $\tfrac12$ in $h$ is the
normalization that conserves cell averages (sums to one), so a constant
field is mapped to itself.

In 3D the same construction applies separately along each axis. The
kernel \emph{factorizes} into a tensor product,
\begin{equation}
  h_{3{\rm D},\,abc} \;=\; h_a\,h_b\,h_c
                     \;=\; \begin{cases}
                              1/8 & a,b,c\in\{0,1\},\\
                              0   & \text{otherwise},
                            \end{cases}
\end{equation}
i.e.\ a $2\!\times\!2\!\times\!2$ array of value $1/8$ (which is
$1/8 = (1/2)^3$, the per-cell weight that makes the kernel sum to
unity). The convolution is
\begin{equation}
  y_{\ell,m,n} \;=\; \sum_{a,b,c} h_{3{\rm D},\,abc}\,
                    \delta_{2N,\,\ell+a,\,m+b,\,n+c},
\end{equation}
a moving average of the eight fine samples in a $2\!\times\!2\!\times\!2$
window. Subsampling $2{:}1$ in each direction then keeps only the
even-indexed entries, so
\begin{equation}
  \delta_{N,IJK} \;=\; y_{2I,2J,2K}
                  \;=\; \frac{1}{8}\!\!\sum_{a,b,c\in\{0,1\}}\!\!
                        \delta_{2N,\,2I+a,\,2J+b,\,2K+c},
\end{equation}
exactly the local mass-conservation formula~(\ref{eq:local-mass}). The
2D version is the same construction with $h_{2{\rm D}}[a,b]=1/4$ on a
$2\!\times\!2$ block. To name the two steps explicitly: forming
$y_{\ell,m,n}$ via the convolution is the \emph{smoothing} step (it
stays on the fine grid and replaces each fine sample by a moving
average of its $2{\times}2{\times}2$ neighbourhood), and keeping only
the entries $y_{2I,2J,2K}$ is the \emph{subsampling} step (it discards
every other sample along each axis to land on the coarse grid).

In the continuous-volume limit (many fine samples per coarse cell, so
$\Delta x_f\to 0$ at fixed $\Delta x_c$) this discrete kernel approaches
a real-space top-hat of width $\Delta x_c$ along each axis (one coarse
cell), followed by point sampling at the coarse-cell centres.

\paragraph{Smoothing step in Fourier space.}
The continuous-$k$ Fourier transform of the discrete kernel $h_n$
(discrete in samples $n$, continuous in $k$ --- not to be confused with
the finite-in/finite-out discrete Fourier transform computed by FFT
routines) is
\begin{equation}
  H(k)\;=\;\sum_n h_n\,e^{-i k n \Delta x_f}
       \;=\;\tfrac12\bigl(1 + e^{-ik\Delta x_f}\bigr)
       \;=\;e^{-ik\Delta x_f/2}\cos\!\left(\tfrac{k\Delta x_f}{2}\right).
\end{equation}
The factored form on the right separates two things. The cosine is the
Fourier transform of the \emph{centered} version of the same kernel
(weights $\tfrac12$ at $x=\pm\Delta x_f/2$), which is real and even
because that kernel is symmetric about the origin. The exponential
$e^{-ik\Delta x_f/2}$ is just the shift property of Fourier transforms
($f(x-s)\leftrightarrow e^{-iks}\hat f(k)$) telling us that our kernel,
sitting at $\{0,\Delta x_f\}$, is offset from the centered one by
$\Delta x_f/2$ --- the center of mass of the two samples we're
averaging.

More generally, the \emph{cell window} $W(\mathbf k)$ is just the
Fourier transform of whatever real-space kernel is used in the
smoothing step of restriction. For \emph{any} averaging operation
written as a convolution
\begin{align}
  y(\mathbf x) &\;=\; \int h(\mathbf x - \mathbf x')\,f(\mathbf x')\,d^d x'
                 \quad\text{(continuous)},\\
  y_{\mathbf m} &\;=\; \sum_{\mathbf n} h_{\mathbf n}\,f_{\mathbf m + \mathbf n}
                 \quad\text{(discrete)},
\end{align}
the cell window is
\begin{equation}
  W(\mathbf k) \;\equiv\; \hat h(\mathbf k),
\end{equation}
where $\hat h$ is the Fourier transform of $h$ (the continuous-$k$
Fourier transform discussed earlier, in the discrete-kernel case). Operationally, $W(\mathbf k)$ is the per-mode
\emph{transfer function} of the smoothing step alone: the smoothed
field's Fourier coefficient at mode $\mathbf k$ is
$\hat y(\mathbf k) = W(\mathbf k)\,\hat f(\mathbf k)$. Different choices
of cell averaging (top-hat, triangular cloud-in-cell, Gaussian, etc.)
produce different $W$, but in every case the same statement holds: $W$
is the Fourier transform of the kernel, and it tells you how strongly
each Fourier mode is multiplied during smoothing. The full restriction
then composes this with the alias sum from subsampling, as in
(\ref{eq:restrict-fourier}).

We adopt the convention of taking the Fourier transform of the
\emph{centered} kernel as our cell window, so that $W$ is real and we
do not have to carry the phase factor around. The phase is just a
half-cell real-space shift; it does not change which Fourier modes are
suppressed by what amount. With this convention,
\begin{equation}\label{eq:Wcell}
  W(\mathbf k) \;=\; \prod_a \cos\!\left(\frac{k_a\Delta x_f}{2}\right)
  \qquad (a\in\{x,y,z\}\text{ in 3D},\; \{x,y\}\text{ in 2D}).
\end{equation}
$W$ is signed real: it can take negative values when one of the
arguments $k_a\Delta x_f/2$ exceeds $\pi/2$, i.e.\ outside the
fine-mode region. Inside the fine-mode region $|k_a|\le\pi/\Delta x_f$,
all arguments lie in $[-\pi/2,\pi/2]$ and $W\ge 0$.

For comparison, the centered \emph{continuous} top-hat of width
$\Delta x_c=2\Delta x_f$ has Fourier transform
\begin{equation}
  W_{\rm cont}(\mathbf k) \;=\; \prod_a \mathrm{sinc}\!\left(\frac{k_a\Delta x_c}{2}\right),
  \qquad \mathrm{sinc}(x)\equiv \sin x/x,
\end{equation}
also signed real (with negative side lobes outside the central peak).
The two definitions are now on equal footing: both are FTs of centered
kernels, no absolute values.

As functions of $\mathbf k$, both $W$ and $W_{\rm cont}$ are defined on
all of $\mathbb R^d$ (the cosines and sincs accept any real argument,
with $\mathrm{sinc}(0)=1$). In our setup the wavenumber is not really
a continuous variable, though: in a periodic box of side $L$ the
allowed modes form the discrete reciprocal lattice $\mathbf k =(2\pi/L)\mathbf m$
with $\mathbf m\in\mathbb Z^d$, and on the fine grid the distinct modes
are further restricted to one period of the lattice, conventionally
$|m_a|\le N$ (equivalently $|k_a|\le\pi/\Delta x_f$). The bare cosine
$\cos(k_a\Delta x_f/2)$ has period $4\pi/\Delta x_f$ in $k_a$, so
inside one fine-mode period $|k_a|\le\pi/\Delta x_f$ it traces out one
half of its cycle, running monotonically from $1$ at $k_a=0$ down to
$0$ at the fine Nyquist. The vanishing of $W$ at the fine Nyquist is
not a coincidence: the fine-Nyquist mode is the alternating pattern
$(+1,-1,+1,-1,\dots)$, and averaging any two adjacent fine samples of
that pattern gives zero, so the smoothing step annihilates this mode
exactly. $W_{\rm cont}$, by contrast, is non-periodic --- it just
decays with side lobes as $|k|\to\infty$.

The periodicity distinction is largely academic for our purposes. In
the periodic-box setup, $\mathbf k$ is not a continuous variable: it
takes only the discrete reciprocal-lattice values
$\mathbf k=(2\pi/L)\mathbf m$ for integer $\mathbf m$, and on a fine
grid of $2N$ points per side the distinct discrete modes are
restricted to one period $|m_a|\le N$, equivalently
$|k_a|\le\pi/\Delta x_f$ (the fine-mode region). So in practice we
never evaluate $W$ outside that fine-mode region, even though $W$ as
a mathematical function is defined on all of $\mathbb R^d$. The
periodicity of $W$ never even gets exercised --- all our evaluations
stay within one fundamental period of the lattice.

$W$ and $W_{\rm cont}$ agree to leading order in $k\Delta x_f$ but
diverge near the coarse \emph{Nyquist wavenumber} (the highest
wavenumber a grid of spacing $\Delta x_c$ can resolve,
$k_{\rm Ny,N} \equiv \pi/\Delta x_c$); in our 2-sample restriction,
equation~(\ref{eq:Wcell}) is the exact one. Figure~\ref{fig:window}
shows the two kernels side by side in real and Fourier space. At the
coarse Nyquist they read $W=\cos(\pi/4)\approx 0.707$ and
$W_{\rm cont}=\mathrm{sinc}(\pi/2)=2/\pi\approx 0.637$, an
$\sim\!11\%$ disagreement.

\begin{figure}[!t]
\centering
\includegraphics[width=\textwidth]{window_compare.pdf}
\caption{Comparison of the centered discrete restriction kernel (blue,
2-sample at $x=\pm\Delta x_f/2$) and the centered continuous top-hat of
width $\Delta x_c$ (red, over $[-\Delta x_c/2,\Delta x_c/2]$). Left:
real-space shapes, with axes in units of $\Delta x_f$. Both have the
same total ``mass'' (sum of stem weights $=$ area under top-hat $=1$)
and the same centroid (the origin). Right: their 1D Fourier-space
windows, $W=\cos(k\Delta x_f/2)$ and
$W_{\rm cont}=\mathrm{sinc}(k\Delta x_c/2)$, plotted on
$u\equiv k\Delta x_f\in[0,2\pi]$ so the side lobes are visible. The
two agree at low $k$, diverge near the coarse Nyquist
($u=\pi/2$, vertical grey line), and both go negative beyond the fine
Nyquist ($u=\pi$, vertical black line).}
\label{fig:window}
\end{figure}

\paragraph{Subsampling step in Fourier space.}
Now the second operation: throw away every other fine sample to land
on the coarse grid. The price is that some fine-grid Fourier modes
become indistinguishable from each other once we only look at every
other point. To see why, evaluate two pure plane waves of fine-grid
wavenumbers $k$ and $k+2\pi/\Delta x_c$ at the coarse-grid positions
$x_I = I\Delta x_c$:
\begin{equation}
  e^{i(k+2\pi/\Delta x_c)\,x_I}
  \;=\; e^{ik\,x_I}\;e^{i2\pi I}
  \;=\; e^{ik\,x_I},
\end{equation}
since $e^{i2\pi I}=1$ for any integer $I$. The two waves give exactly
the same sequence of values on the coarse grid, so the coarse grid
cannot tell them apart --- on the coarse grid, the higher-frequency
mode is indistinguishable from the lower-frequency one. This is what
``\emph{aliasing}'' means: a high-frequency mode that the coarse grid
cannot resolve gets misidentified as a lower-frequency mode it can. The shift between aliased modes is
$\Delta k_{\rm alias}=2\pi/\Delta x_c=\pi/\Delta x_f$ (recall
$\Delta x_c=2\Delta x_f$). Note that this alias spacing along each
axis is exactly the fine-grid Nyquist wavenumber $k_{\rm Ny,2N}$:
every coarse-grid mode at $k_c$ pairs with the fine-grid mode at
$k_c + k_{\rm Ny,2N}$, which sits in the upper half of the fine-mode
region (between coarse and fine Nyquist), and the two are
indistinguishable after subsampling. Figure~\ref{fig:alias} shows the
2D version of this pairing: in 2D, each coarse mode receives
contributions from $2\times 2 = 4$ fine modes (one principal and three
aliases), one in each quadrant of the fine-mode region. In 3D the
count becomes $2\times 2\times 2 = 8$: one principal and seven
aliases, one in each of the seven other octants.

\begin{figure}[!t]
\centering
\includegraphics[width=0.7\textwidth]{alias_diagram.pdf}
\caption{Alias pairing under $2{:}1$ subsampling in 2D. The fine-mode
region (light red square) splits into a central coarse zone (light blue,
$|k_a|\le k_{\rm Ny,N}$) and three outer-wing regions. A representative
coarse mode $k_c$ in the coarse zone (filled blue circle) receives
contributions from itself (the principal $\mathbf n=(0,0)$ term) and
from three alias partners at $k_c + \pi\mathbf n/\Delta x_f$ for
$\mathbf n\in\{(1,0),(0,1),(1,1)\}$. The arrows show the folding: each
alias partner, after periodic identification, sits in one of the three
other quadrants of the fine-mode region.}
\label{fig:alias}
\end{figure}

Concretely, for every coarse mode $\mathbf k_c$ inside the resolved
coarse-mode region ($|k_{c,a}|\le\pi/\Delta x_c$ for each axis), there
are exactly two fine-grid modes per axis that fall inside the fine-mode
region and alias onto it: one at $k_{c,a}$ itself (the ``principal''
mode, already in the coarse region) and one at
$k_{c,a}\pm\pi/\Delta x_f$ (with the sign chosen so that the result
lies inside the fine-mode region; for positive $k_{c,a}$ this means
$k_{c,a}-\pi/\Delta x_f$, in the negative-side outer wing of the fine
zone). Either way, the alias partner has magnitude
$|k_a|\in[k_{\rm Ny,N},\,k_{\rm Ny,2N}]$, sitting beyond the coarse
Nyquist but inside the fine Nyquist. In 3D this gives $2\times 2\times 2 = 8$
fine modes per coarse mode (one principal + seven aliases),
enumerated by which axis is shifted or not.
Calling that bookkeeping vector $\mathbf n\in\{0,1\}^3$:
\begin{equation}\label{eq:restrict-fourier}
  \hat\delta_N(\mathbf k_c)
  \;=\; \sum_{\mathbf n\in\{0,1\}^3}
        W(\mathbf k_c + \pi\mathbf n/\Delta x_f)\,
        \hat\delta_{2N}(\mathbf k_c + \pi\mathbf n/\Delta x_f).
\end{equation}
The cell window $W$ from~(\ref{eq:Wcell}) multiplies each contribution
because the convolution step happened first, before subsampling. The
$\mathbf n=\mathbf 0$ term is the \emph{principal} contribution: it is
the cell-window-suppressed value of $\hat\delta_{2N}$ at $\mathbf k_c$
itself. The seven $\mathbf n\ne\mathbf 0$ terms are the
\emph{aliases}: they carry fine-grid power from \emph{outside} the
coarse Nyquist back into the coarse mode $\mathbf k_c$, where it is no
longer separable from the genuine low-frequency content. Appendix
\ref{app:aliassum} gives a derivation of (\ref{eq:restrict-fourier})
that tracks the Fourier-convention prefactor explicitly.

The two pieces play complementary roles. The convolution
(\ref{eq:Wcell}) acts as a smoothing filter: it suppresses high-$k$
power inside the resolved coarse-mode region smoothly, with $W=1$ at
$\mathbf k=\mathbf 0$ and $W\!\to\!0$ at the corners of the fine-mode
region. The subsampling step keeps every other fine sample in each
direction; equivalently, in real space it multiplies by a sum of
equally-spaced spikes that pick out only the kept samples. In Fourier
space, this folds the fine spectrum onto the coarse one. Picture the
fine spectrum living inside the fine-mode region $|k_a|\le\pi/\Delta x_f$;
after subsampling, the new fundamental period along each axis is half
as wide, $|k_a|\le\pi/\Delta x_c$. To get the new spectrum, you overlay
the fine spectrum on itself, shifted by every integer multiple of the
new period $2\pi/\Delta x_c=\pi/\Delta x_f$ along each axis, and sum
within one new fundamental cell. Inside the larger fine-mode region
there are two ``shift or no shift'' choices per axis, so in 3D each
coarse-grid mode receives $2\times 2\times 2 = 8$ overlapping
contributions --- one principal ($\mathbf n=\mathbf 0$, the unshifted
copy) and seven aliases (shifted copies). These are exactly the eight
terms in (\ref{eq:restrict-fourier}). Restriction is therefore not a clean
band-limiting operation --- it is a smoothing step followed by a fold.
In particular, restriction is \emph{not} simply ``multiply by $W$ in
Fourier space.'' Multiplication by $W$ is only the smoothing step,
which produces the principal ($\mathbf n=\mathbf 0$) term in
(\ref{eq:restrict-fourier}); the seven alias terms come from the
subsampling. The two pictures coincide only in the special case where
$\hat\delta_{2N}$ has no power above the coarse Nyquist
($|k_a|>\pi/\Delta x_c$), so that all $\mathbf n\ne\mathbf 0$ terms
vanish identically.

The same alias structure shows up at the level of power spectra. For
the smoothing step alone, smoothing-as-convolution becomes
$\hat y(\mathbf k) = W(\mathbf k)\,\hat\delta_{2N}(\mathbf k)$ in
Fourier space, so the smoothed-field power spectrum is
$P_y(\mathbf k) = |W(\mathbf k)|^2\,P_{2N}(\mathbf k)$ --- the standard
``convolution multiplies $P$ by $|W|^2$'' result. After the
subsampling step, taking the squared modulus of
(\ref{eq:restrict-fourier}) and using the fact that distinct Fourier
modes of a statistically homogeneous Gaussian field are uncorrelated
(so cross-terms in $\mathbf n\ne\mathbf n'$ average to zero), the
\emph{coarse-grid} power spectrum becomes the alias sum
\begin{equation}\label{eq:Pk-alias}
  P_N(\mathbf k_c) \;=\; \sum_{\mathbf n\in\{0,1\}^3}
        \bigl|W(\mathbf k_c+\pi\mathbf n/\Delta x_f)\bigr|^2\,
        P_{2N}(\mathbf k_c+\pi\mathbf n/\Delta x_f).
\end{equation}
The $\mathbf n=\mathbf 0$ term is the clean ``$|W|^2 P$'' smoothing
contribution; the seven $\mathbf n\ne\mathbf 0$ terms add fine-grid
power from beyond the coarse Nyquist, weighted by $|W|^2$ evaluated at
those higher wavenumbers --- a positive, incoherent contamination of
the coarse spectrum.

Figure~\ref{fig:density} compares $\delta_{2N}$ and $\delta_N$ in real
space (panels~1 and~2) and shows the Fourier amplitude map of the fine
field (panel~3). Panel~4 plots the per-mode amplitude ratio
$|\hat\delta_N(\mathbf k_c)|/|\hat\delta_{2N}(\mathbf k_c)|$ against
$|\mathbf k_c|$ on a linear-$k$, log-$|$ratio$|$ scale. Three things
are overlaid: the measured ratios (blue), the exact 2D cell window
$W^2(k_x,k_y)$ evaluated at the same modes (green), and the 1D proxy
$\cos^2(k\Delta x_f/2)$ (red curve). The blue and green points
together expose the two ingredients of the alias-sum formula: the
green points show what restriction would produce if there were no
aliasing (just the principal-mode multiplication by $W$), and the blue
points include the alias contributions on top.

Note that $W$ is direction-dependent: at fixed $|\mathbf k|$, modes
along an axis and modes along the diagonal feel different cell-window
suppression. For example at the coarse Nyquist along an axis,
$\mathbf k_c=(\pi/\Delta x_c, 0)$, the principal suppression is
$W=\cos(\pi/4)\cdot 1 \approx 0.707$, but at the corner of the coarse
zone $\mathbf k_c=(\pi/\Delta x_c, \pi/\Delta x_c)$ it drops to
$W=\cos^2(\pi/4)=1/2$ (and in 3D the cube corner gives
$\cos^3(\pi/4)\approx 0.354$). At each $|\mathbf k|$ the green points
therefore form a narrow band rather than a single value, with the 1D
proxy curve lying somewhere inside that band. The much larger blue
scatter on top of the green band is the alias contamination, which
grows toward the coarse Nyquist where high-$k$ fine-grid power
($|\mathbf k|>\pi/\Delta x_c$) is folded back in by subsampling.

\begin{figure}[!t]
\centering
\includegraphics[width=\textwidth]{restriction_density.pdf}
\caption{Restriction of $\delta_{2N}$ to $\delta_N$ by $2\times2$ block
averaging. Top row: real-space maps of $\delta_{2N}$, $\delta_N$, and
$\log_{10}|\hat\delta_{2N}|$. Bottom row: per-mode amplitude ratio
(blue) overlaid with the exact 2D cell window $W^2(k_x,k_y)$ (green)
and the 1D proxy $\cos^2(k\Delta x_f/2)$ (red); and the radially
binned power ratio $P_N(k_c)/P_{2N}(k_c)$ (blue circles) compared with
the alias-sum theory prediction (red squares). The two curves should
agree if the alias-sum formula~(\ref{eq:Pk-alias}) is correct.}
\label{fig:density}
\end{figure}

\section{1LPT displacement under restriction}

\emph{Lagrangian perturbation theory} (LPT) describes structure
formation by following individual fluid elements through their
displacement from an initial uniform configuration. Each fluid element
is labelled by its initial (Lagrangian) position $\mathbf q$, and its
later (Eulerian) position is $\mathbf x(\mathbf q, t) = \mathbf q +
\boldsymbol\Psi(\mathbf q, t)$, where $\boldsymbol\Psi$ is the
displacement. To first order in the density contrast $\delta$, the
displacement $\boldsymbol{\Psi}^{(1)}$ satisfies
\begin{equation}\label{eq:lpt-poisson}
  \nabla_q \cdot \boldsymbol{\Psi}^{(1)}(\mathbf q) = -\delta_{\rm lin}(\mathbf q),
\end{equation}
which is mass conservation kept only to first order in
$\delta$ (the full nonlinear continuity equation reduces to this when
products of small perturbations are dropped). $\delta_{\rm lin}$ here
is the linear-theory density contrast at some reference time. (See
\texttt{notes/cosmo\_ic.tex} for the full derivation and the
extension to higher orders.)

The growing-mode solution to (\ref{eq:lpt-poisson}) --- the one whose
amplitude grows in time as the universe evolves, as opposed to the
decaying mode --- is curl-free, so we can write
$\boldsymbol{\Psi}^{(1)} = -\nabla_q\phi^{(1)}$ for some scalar
potential $\phi^{(1)}$. Substituting into~(\ref{eq:lpt-poisson})
yields a Poisson-like equation
$\nabla_q^2\phi^{(1)} = \delta_{\rm lin}$ for the
\emph{displacement potential} (Laplacian with respect to the Lagrangian
coordinate $\mathbf q$). The scalar $\phi^{(1)}$ is also called the
Lagrangian potential; it is distinct from the gravitational potential
$\Phi$ (which obeys $\nabla^2\Phi = 4\pi G\bar\rho a^2\delta$ and
carries different units), though the two are proportional in linear
theory. They solve Poisson-type equations with structurally similar
but physically different right-hand sides.

In Fourier space the Laplacian becomes multiplication by $-k^2$, so the
Poisson equation $\nabla_q^2\phi^{(1)} = \delta_{\rm lin}$ gives
\begin{equation}
  -k^2\,\hat\phi^{(1)}(\mathbf k) \;=\; \hat\delta_{\rm lin}(\mathbf k)
  \;\Longrightarrow\;
  \hat\phi^{(1)}(\mathbf k) \;=\; -\frac{\hat\delta_{\rm lin}(\mathbf k)}{k^2}.
\end{equation}
Substituting into $\boldsymbol{\Psi}^{(1)} = -\nabla_q\phi^{(1)}$
(gradient becomes $i\mathbf k$ in Fourier),
\begin{equation}\label{eq:psi-kernel}
  \hat{\boldsymbol{\Psi}}^{(1)}(\mathbf k)
  \;=\; -i\mathbf k\,\hat\phi^{(1)}(\mathbf k)
  \;=\; -i\mathbf k\cdot\left(-\frac{\hat\delta_{\rm lin}(\mathbf k)}{k^2}\right)
  \;=\; +\frac{i\mathbf k}{k^2}\,\hat\delta_{\rm lin}(\mathbf k).
\end{equation}
As a sanity check, the Poisson constraint is satisfied identically:
$i\mathbf k\!\cdot\!\hat{\boldsymbol\Psi}^{(1)} = (i\mathbf k)\!\cdot\!(i\mathbf k\hat\delta/k^2) = -\hat\delta_{\rm lin}$, in agreement with~(\ref{eq:lpt-poisson}).

The peculiar velocity (the deviation from Hubble flow) follows from
$\mathbf v^{(1)} = aHf\,\boldsymbol{\Psi}^{(1)}$, where $a$ is the
cosmological scale factor, $H=\dot a/a$ is the Hubble rate, and
$f\equiv d\ln D/d\ln a$ is the linear growth rate ($D$ being the
linear growth factor). The kernel $i\mathbf k/k^2$ in
(\ref{eq:psi-kernel}) weights \emph{large scales} (the $1/k^2$
suppression dominates the linear $k$ in the numerator), so restriction
has only a mild effect on $\boldsymbol{\Psi}^{(1)}$.

Figure~\ref{fig:psi} compares two coarse-grid versions of $\Psi_x^{(1)}$:
\begin{itemize}
\item[(a)] solve \eqref{eq:psi-kernel} on the fine grid, then block-average
      each component to the coarse grid;
\item[(b)] restrict $\delta$ first, then solve \eqref{eq:psi-kernel} on
      the coarse grid.
\end{itemize}
The two are visually indistinguishable on the colour scale of the field
itself, and the difference panel shows a residual that is roughly two
orders of magnitude smaller. The kernel $i\mathbf k/k^2$ is
large-scale-dominated, which is why 1LPT is robust to restriction.

\begin{figure}[!t]
\centering
\includegraphics[width=\textwidth]{restriction_psi.pdf}
\caption{1LPT displacement $\Psi_x^{(1)}$ on the coarse grid. Left:
fine-grid $\Psi^{(1)}$ block-averaged. Middle: $\Psi^{(1)}$ obtained
from the restricted $\delta_N$. Right: their difference (note the
unchanged colour scale).}
\label{fig:psi}
\end{figure}

\section{Deformation tensor and the 2LPT source}\label{sec:deformation}

The gradient of the displacement, $\Psi^{(1)}_{i,j} \equiv \partial_i\Psi^{(1)}_j$
(the comma in the subscript denotes a partial derivative, $f_{,i}\equiv\partial f/\partial q_i$),
is the Lagrangian \emph{deformation tensor}. Its trace is the Poisson
constraint $\sum_i\Psi^{(1)}_{i,i} = -\delta_{\rm lin}$, and the velocity gradient
follows from $\partial_i v^{(1)}_j = aHf\,\Psi^{(1)}_{i,j}$. More importantly
for our purposes, the deformation tensor is the building block of the
2LPT source $S^{(2)}$, which appears as the right-hand side of a
second Poisson-like equation
$\nabla_q^2\phi^{(2)} = S^{(2)}$ for the second-order displacement
potential $\phi^{(2)}$ (with $\boldsymbol\Psi^{(2)} = -\nabla_q\phi^{(2)}$).
In two dimensions the analogue of the 3D 2LPT source has the single
off-diagonal invariant
\begin{equation}\label{eq:s2}
  S^{(2)}(\mathbf q) \;=\; \Psi^{(1)}_{x,x}\,\Psi^{(1)}_{y,y}
                          \;-\;\bigl(\Psi^{(1)}_{x,y}\bigr)^2,
\end{equation}
and in 3D $S^{(2)}=\sum_{i<j}\!\bigl[\Psi^{(1)}_{i,i}\Psi^{(1)}_{j,j}-(\Psi^{(1)}_{i,j})^2\bigr]$.
Because $S^{(2)}$ is quadratic in $\Psi^{(1)}_{i,j}$, any restriction
error in the deformation tensor enters $S^{(2)}$ at second order: a
small fractional error in $\Psi^{(1)}_{i,j}$ produces a fractional error
of comparable size in each of the two products in (\ref{eq:s2}), so the
2LPT source is the most sensitive LPT object to restriction.

The Fourier kernel of $\Psi^{(1)}_{i,j}$ is
\begin{equation}
  \widehat{\Psi^{(1)}_{i,j}}(\mathbf k) \;=\; -\frac{k_i k_j}{k^2}\,\hat\delta(\mathbf k),
\end{equation}
whose magnitude is bounded by unity --- the deformation tensor neither
amplifies nor strongly suppresses any scale. So unlike either
$\nabla\delta$ or $\boldsymbol{\Psi}^{(1)}$, the deformation tensor does
not have an obvious bias towards small or large scales. Restriction
errors come from the cell window $W$ multiplying the kernel, plus alias
contributions of the same form.

Figure~\ref{fig:deformation} compares (top row) the diagonal component
$\Psi^{(1)}_{x,x}$ obtained from the fine-then-restrict path against the
restrict-then-solve path, and (bottom row) the 2LPT source $S^{(2)}$
constructed both ways. The $\Psi$-component error is small and concentrated
where the field has fine structure; the $S^{(2)}$ error is visibly
larger, consistent with the quadratic propagation. This is the
practically interesting cost of restriction in an LPT pipeline: 1LPT
displacements survive well, but 2LPT sources degrade faster.

\begin{figure}[!t]
\centering
\includegraphics[width=\textwidth]{restriction_deformation.pdf}
\caption{Deformation tensor and 2LPT source under restriction.
Top row: $\Psi^{(1)}_{x,x}=\partial_x\Psi_x^{(1)}$ from the
fine-then-restrict path, from the restrict-then-solve path, and their
difference. Bottom row: the 2D 2LPT source $S^{(2)} = \Psi^{(1)}_{x,x}\Psi^{(1)}_{y,y} - (\Psi^{(1)}_{x,y})^2$
constructed both ways and their difference. The $S^{(2)}$ error is
larger because restriction errors in the deformation tensor enter $S^{(2)}$ quadratically.}
\label{fig:deformation}
\end{figure}

\section{Restricting the displacement potential is Poisson-inconsistent}\label{sec:phi}

A natural-looking shortcut is: solve the Poisson equation once on the
fine grid for $\phi^{(1)}_{2N}$, restrict $\phi^{(1)}_{2N}\to\phi^{(1)}_N$,
then take the gradient on the coarse grid. This is mathematically distinct from
restricting $\delta$ and then solving Poisson. Three coarse-grid paths
for $\Psi_x^{(1)}$ are summarised below. Writing the contribution of a
fine-grid mode $\mathbf k'=\mathbf k_c+2\pi\mathbf n/\Delta x_c$ to the
coarse mode $\mathbf k_c$ as a prefactor times
$\hat\delta_{2N}(\mathbf k')$:
\begin{center}
\begin{tabular}{lll}
\hline
path & principal ($\mathbf n=\mathbf 0$) & alias ($\mathbf n\ne\mathbf 0$) \\
\hline
(1) fine $\Psi$, then restrict
   & $i k_{c,x}\,W(\mathbf k_c)/k_c^{\,2}$
   & $i k'_x\,W(\mathbf k')/(k')^{2}$ \\
(2) restrict $\delta$, then solve
   & $i k_{c,x}\,W(\mathbf k_c)/k_c^{\,2}$
   & $i k_{c,x}\,W(\mathbf k')/k_c^{\,2}$ \\
(3) restrict $\phi^{(1)}$, then $-\nabla\phi^{(1)}$
   & $i k_{c,x}\,W(\mathbf k_c)/k_c^{\,2}$
   & $i k_{c,x}\,W(\mathbf k')/(k')^{2}$ \\
\hline
\end{tabular}
\end{center}
The principal terms agree. The alias contributions show that path~(2)
amplifies aliased high-$k$ power by $1/k_c^{\,2}$, while path~(3) keeps
the original $1/(k')^{2}$ suppression and is in that sense closer to
path~(1) in amplitude.

That, however, is not the right diagnostic. The deeper problem with
path~(3) is \emph{Poisson inconsistency}. The restricted potential
$\phi^{(1)}_N=R\,\phi^{(1)}_{2N}$ does not solve the coarse-grid
Poisson equation sourced by the restricted density:
\begin{equation}
  -\nabla^2 \phi^{(1)}_N \;\ne\; \delta_N.
\end{equation}
Equivalently, the effective density implied by $\phi^{(1)}_N$ via the
coarse-grid Laplacian,
\begin{equation}
  \delta_{\rm eff}(\mathbf q) \;\equiv\; -\nabla^2\phi^{(1)}_N(\mathbf q),
\end{equation}
is not equal to $\delta_N$. The displacement obtained as
$-\nabla\phi^{(1)}_N$ is therefore not the 1LPT displacement of
$\delta_N$ --- it is the 1LPT displacement of some other field
$\delta_{\rm eff}$.

Figure~\ref{fig:phi} shows the three paths and the
$\delta_{\rm eff}-\delta_N$ residual. The top row shows the three $\Psi_x$
maps; visually all three look similar, consistent with the analytic
table above. The bottom-left and bottom-middle panels are the error
maps against the fine reference. The bottom-right panel is the
self-consistency residual: it is structured (not noise), with a
clear correlation to the underlying density field, demonstrating that
$\phi^{(1)}_N$ encodes a different source than $\delta_N$.

\begin{figure}[!t]
\centering
\includegraphics[width=\textwidth]{restriction_phi.pdf}
\caption{Three coarse-grid paths to $\Psi_x^{(1)}$ and the Poisson
inconsistency of restricting $\phi^{(1)}$. Top: the three $\Psi_x$ maps.
Bottom: error vs the fine reference for paths~(2) and~(3), and
$-\nabla^2 R\phi^{(1)}_{2N}-\delta_N$ (the restricted potential does
not solve the coarse Poisson equation sourced by the restricted
density).}
\label{fig:phi}
\end{figure}

\section{Comparison with a CDM-like spectrum}\label{sec:cdm}

The toy $P(k)\propto k^n$ used above (with $n=-2$ in our figures) is
chosen as the 2D analog of the high-$k$ behavior of the physical
$\Lambda$CDM (cold dark matter with cosmological constant) matter
power spectrum. In 3D, $\Lambda$CDM at high $k$ (well above the
matter-radiation equality wavenumber
$k_{\rm eq}\sim 0.01\,h\,\text{Mpc}^{-1}$) behaves as
$P_{3D}(k)\sim k^{n_s-4}\sim k^{-3}$, where $n_s\approx 0.965$ is the
primordial spectral tilt (with a logarithmic correction from the
radiation-era transfer function).

The literal slope $n=-3$ is a 3D quantity and does not translate
directly to a 2D demonstration. The relevant invariant for restriction
effects is the variance contribution per logarithmic $k$ bin:
\begin{equation}
  \frac{d\sigma^2}{d\ln k} \;\propto\;
  \begin{cases}
    k^{3}\,P_{3D}(k) & \text{(3D)},\\[2pt]
    k^{2}\,P_{2D}(k) & \text{(2D)}.
  \end{cases}
\end{equation}
For 3D $\Lambda$CDM at high $k$, $k^3 P_{3D}\sim k^0$ --- a constant
contribution per log bin (so the variance integral diverges only
logarithmically). Imposing the same flat-per-log-bin behavior in 2D
requires $k^2 P_{2D}\sim k^0$, i.e.\ $P_{2D}\propto k^{-2}$. That is
the 2D power spectrum whose window and alias effects under restriction
behave qualitatively like those of 3D $\Lambda$CDM, and that is why
we use $n=-2$ in the figures.

The toy still misses the low-$k$ turnover at $k_{\rm eq}$ where the
$\Lambda$CDM spectrum bends from $P(k)\sim k$ (large scales) to
$P(k)\sim k^{-3}$ (small scales). To check whether the conclusions
change qualitatively when the input includes that turnover,
Figure~\ref{fig:cdm} repeats the comparisons using a $\Lambda$CDM
matter power spectrum at $z=200$, generated with \texttt{CLASS} (a
public Boltzmann code that integrates the linear cosmological
perturbation equations through the radiation, matter, and dark-energy
eras), on a 2D grid with $L=500\,h^{-1}$~Mpc, $2N=128$, $N=64$. (The
length unit $h^{-1}\,\mathrm{Mpc}$ is the standard cosmology
convention: $h$ is the dimensionless Hubble parameter, $H_0 = 100h\,
\mathrm{km}\,\mathrm{s}^{-1}\,\mathrm{Mpc}^{-1}$. See
\texttt{notes/cosmo\_ic.tex} for cosmology background.)

The top-left panel shows the input: the true CDM spectrum (black) versus
the toy $k^n$ used in the previous figures (red dashed), with the
coarse and fine Nyquist wavenumbers marked. The CDM spectrum has the
matter-radiation-equality turnover at low $k$ that the toy lacks, but
the high-$k$ slopes are similar. The top-right
two panels show $\delta_{2N}$ and $\delta_N$; the bottom row shows the
residuals for the mass-conservation sanity check ($R\delta_{2N}-\delta_N=0$),
the 1LPT displacement $\Psi_x^{(1)}$, and the 2LPT source $S^{(2)}$,
each on a colour scale set to a fraction of the field's own dynamic range.

The qualitative conclusions are unchanged: $\Psi^{(1)}$ is highly robust,
$S^{(2)}$ shows visible residuals, and the Poisson-inconsistency story
of \S\ref{sec:phi} is independent of the choice of input spectrum. The
absolute relative errors are smaller with the CDM spectrum (printed at
the bottom of the figure script), simply because there is less power for
the cell window and aliasing to act on near the coarse Nyquist.

\begin{figure}[!t]
\centering
\includegraphics[width=\textwidth]{restriction_cdm.pdf}
\caption{Restriction with a realistic CDM-like spectrum (CLASS, $z=200$,
$L=500\,h^{-1}$~Mpc). Top-left: input $P(k)$ (black) versus the toy
$k^n$ used in the earlier figures (red), with coarse/fine Nyquist
wavenumbers marked. Top-middle/right: $\delta_{2N}$ and $\delta_N$.
Bottom row, left to right: mass-conservation sanity ($R\delta_{2N}-\delta_N$,
identically zero), $\Psi_x^{(1)}$ residual, and $S^{(2)}$ residual
between the two coarse-grid construction paths.}
\label{fig:cdm}
\end{figure}

\section{Application: zoom-in IC generation}\label{sec:zoom}

The bin-averaging analysis of \S\ref{sec:fourier} and the
Poisson-inconsistency result of \S\ref{sec:phi} are tools, not
applications. The natural place they get applied is in zoom-in
initial-condition construction, where one wants a fine-resolution
``zoom region'' inside a larger lo-res box.

A naive proposal: generate the full density (or potential) field on the
fine grid everywhere, then construct the lo-res region outside the
zoom by simple block averaging. Two variants of this idea:
\begin{description}
\item[A1: bin-average $\delta$.]
  Set $\delta_{\rm lo} = R\,\delta_{2N}$ outside the zoom, then solve
  Poisson on the coarse grid sourced by $\delta_{\rm lo}$ to get
  $\boldsymbol{\Psi}_{\rm lo}$.
\item[A2: bin-average $\phi^{(1)}$.]
  Compute $\phi^{(1)}_{2N}$ on the fine grid, restrict
  $\phi^{(1)}_{\rm lo} = R\,\phi^{(1)}_{2N}$, and take the gradient on
  the coarse grid: $\boldsymbol{\Psi}_{\rm lo} = -\nabla\phi^{(1)}_{\rm lo}$.
\end{description}
Both fail, in different ways. The standard production approach,
\begin{description}
\item[B: k-truncate $\mu$.]
  Take the same fine-grid white noise $\mu$, zero its modes above the
  coarse Nyquist to get $\mu_{\rm trunc}$, then convolve with the
  transfer function to get
  $\delta_{\rm lo} = \mathcal F^{-1}[T(k)\,\hat\mu_{\rm trunc}(\mathbf k)]$.
  Solve Poisson on the coarse grid sourced by $\delta_{\rm lo}$.
\end{description}
gets it right.

\paragraph{Why A1 fails (cell-window deficit).}
The alias-sum formula~(\ref{eq:restrict-fourier}) gives
\begin{equation}
  R\,\delta_{2N}(\mathbf k_c) \;=\; W(\mathbf k_c)\,\delta_{2N}(\mathbf k_c)
  \;+\; \sum_{\mathbf n\ne\mathbf 0} W(\mathbf k_c+\pi\mathbf n/\Delta x_f)\,\delta_{2N}(\mathbf k_c+\pi\mathbf n/\Delta x_f).
\end{equation}
The principal $W(\mathbf k_c)\,\delta_{2N}(\mathbf k_c)$ term is
suppressed by $W = \cos(k_a\Delta x_f/2)$ along each axis, dropping
to $\cos(\pi/4)\approx 0.71$ at the coarse Nyquist along an axis and
to $\cos^2(\pi/4)=1/2$ at the 2D corner. Inside the zoom region the
field carries all fine modes; outside it has the same low-$k$ modes
suppressed by $W$. The mid-$k$ amplitude mismatch is a
\emph{boundary discontinuity in $\delta$} at the zoom edge.

The lo-res displacement $\boldsymbol{\Psi}_{\rm lo}$ is internally
Poisson-consistent with the lo-res $\delta_{\rm lo}$ (we solved Poisson
on the coarse grid sourced by it), so there is no inconsistency
\emph{within} the lo-res region. The failure is the boundary mismatch
with the hi-res zoom interior.

\paragraph{Why A2 fails (Poisson inconsistency).}
The Poisson-inconsistency result of \S\ref{sec:phi} says
$-\nabla^2(R\,\phi^{(1)}_{2N}) \neq R\,\delta_{2N}$. The displacement
$\boldsymbol{\Psi}_{\rm lo} = -\nabla(R\,\phi^{(1)}_{2N})$ is the
1LPT displacement of an effective density
$\delta_{\rm eff} \equiv -\nabla^2(R\,\phi^{(1)}_{2N})$, which equals
neither the bin-averaged density nor the original $\delta_{2N}$.

Mode by mode, $\delta_{\rm eff}(\mathbf k_c)$ has the same principal
$W(\mathbf k_c)\delta_{2N}(\mathbf k_c)$ contribution as A1 but its
alias contributions are extra-suppressed by $k_c^2/|\mathbf k_c+\pi\mathbf n/\Delta x_f|^2$.
So the cell-window deficit at the boundary is similar to A1's, but
the deeper failure is conceptual: the lo-res $\boldsymbol{\Psi}$ and
$\delta_{\rm lo}$ refer to different fields. Inside the zoom,
particles are displaced consistently with $\delta_{2N}$; outside,
they're displaced consistently with $\delta_{\rm eff}$. The two
regions don't represent the same underlying density distribution.

\paragraph{Why B works (white noise is flat).}
For approach B, the lo-res field is $\delta_{\rm lo}(\mathbf k_c) =
T(\mathbf k_c)\,\hat\mu(\mathbf k_c)$ for $|k_a|\le\pi/\Delta x_c$,
zero otherwise. There is no cell-window suppression and no aliasing,
because we band-limited $\mu$ \emph{before} convolving with $T(k)$.

The reason this works is that white noise has flat power spectrum:
$P_\mu(\mathbf k)=$ constant. In the alias-sum formula
(\ref{eq:Pk-alias}) applied to $\mu$ instead of $\delta$, the
aliased contributions have the same per-mode amplitude as the
principal, so cutting them out (rather than letting them fold back
in) introduces no statistical bias. The truncated white noise is
still white --- just on a smaller fundamental period. Convolving
with $T(k)$ then produces a density field that exactly matches what
a fully resolved hi-res run would give at all $|\mathbf k|\le k_{\rm Ny,N}$.

(In production codes such as \texttt{MUSIC2} v2, the same property is
achieved by generating white noise independently on each level and
splicing in Fourier space: each level provides modes between its
parent's Nyquist and its own. Approach B is a simplified pedagogical
form of the same idea. See \texttt{notes/cosmo\_ic.tex} for the
\texttt{MUSIC2} pipeline in detail and \texttt{notes/ic\_sampling.tex}
for the broader ``$P$-sampled vs $\xi$-sampled'' discussion of how
white noise gets generated in IC codes.)

\paragraph{Numerical demonstration.}
Figure~\ref{fig:zoom-ic-cmp} compares the three approaches on a 2D
toy IC ($L=1$, $N_{\rm fine}=128$, $N_{\rm coarse}=64$, $P(k)\propto k^{-2}$,
hard-step zoom region $[L/4,3L/4]^2$). The diagnostics arranged across
five rows:
\begin{enumerate}
\item Lo-res $\delta_{\rm lo}$ as built (different per approach, see
      caption);
\item Residual against the ideal lowpass
      $\delta_{\rm lo} - \delta_{\rm ideal\,lp}$, where the reference
      is $\mathcal F^{-1}[T(k)\hat\mu_{\rm trunc}]$ from B itself;
\item Poisson inconsistency $-\nabla^2\phi^{(1)}_{\rm lo} - \delta_{\rm lo}$;
\item Composite IC density (hi-res inside zoom, lo-res outside);
\item 1D coarse-grid slice through $y=L/2$, with the ideal lowpass
      overlaid.
\end{enumerate}

Quantitatively (RMS values):
\begin{itemize}
\item \emph{Residual vs ideal lowpass} (row 2): A1 and A2 are
      comparable ($\sim 2\times 10^{-3}$, dominated by the shared
      $W-1$ principal term); B is exactly zero.
\item \emph{Poisson inconsistency} (row 3): A1 and B are at machine
      zero; A2 is $\sim 5\times 10^{-4}$, with structured residual.
      This is what distinguishes A2 from A1: same boundary-deficit
      symptom, but additionally Poisson-inconsistent.
\end{itemize}
Approach B is the only one that produces a lo-res region matching
the ideal band-limited version of $\delta_{2N}$ and that is
internally Poisson-consistent.

\begin{figure}[!t]
\centering
\includegraphics[width=\textwidth]{zoom_ic_comparison.pdf}
\caption{Three approaches to constructing the lo-res region of a
zoom-in IC, compared on a 2D toy problem
($L=1$, $N_{\rm fine}=128$, $N_{\rm coarse}=64$, $P(k)\propto k^{-2}$,
seed 42). Columns: A1 (bin-average $\delta$), A2 (bin-average
$\phi^{(1)}$), B (k-truncate $\mu$, the standard approach). Rows
(top to bottom): lo-res $\delta_{\rm lo}$ as built; residual against
the ideal lowpass; Poisson-inconsistency residual
$-\nabla^2\phi^{(1)}_{\rm lo}-\delta_{\rm lo}$; composite IC density
(hi-res inside the green box, lo-res outside); 1D coarse-grid slice
at $y=L/2$ with the ideal lowpass overlaid.
The per-approach definition of $\delta_{\rm lo}$ in row 1:
A1 uses $R\,\delta_{2N}$; A2 uses
$\delta_{\rm eff}\equiv -\nabla^2(R\,\phi^{(1)}_{2N})$ (the density
implied by the bin-averaged potential, since
$R\,\phi^{(1)}_{2N}$ is not Poisson-consistent with any $R\,\delta$);
B uses the band-limited $\delta$ sampled on the coarse grid. The
ideal-lowpass reference for row 2 is
$\mathcal F^{-1}[T(k)\hat\mu_{\rm trunc}(\mathbf k)]$, sampled on the
coarse grid.}
\label{fig:zoom-ic-cmp}
\end{figure}

\paragraph{Spectral diagnostics.}
Figure~\ref{fig:zoom-ic-pk} restates the same comparison at the level
of binned power spectra. The top row shows the lo-res $P(k)$ of each
restriction method and its ratio to the input theory: A1 and A2 both
develop a deficit growing toward the coarse Nyquist
($|\mathbf k|/k_{\rm Ny,N}\to 1$); B sits flat at the theory line.
For visualisation we use the ``fix-and-pair'' trick \citep{Angulo2016}:
each Fourier mode of the white noise is rescaled to have
$|\hat\mu_{\mathbf k}| = \sqrt{N_{\rm tot}}$ before convolving with
$T(k)$. This removes the per-mode cosmic-variance scatter and exposes
the systematic restriction effects cleanly.

The bottom row of Figure~\ref{fig:zoom-ic-pk} is an aside: it shows
the standard particle-based IC validation. Particles are placed at
fine-grid Lagrangian positions, displaced via 2LPT, and deposited via
2D cloud-in-cell (CIC) onto the fine grid. The raw measured $P(k)$
is suppressed by the CIC window
$W_{\rm CIC}(\mathbf k) = \prod_a\mathrm{sinc}(k_a\Delta x_f/2)$
along each axis (squared in the spectrum); deconvolving by $W_{\rm CIC}^2$
recovers the input theory across all $|\mathbf k|<k_{\rm Ny,2N}$. The
two rows live on different grids (top: coarse; bottom: fine) and
answer different questions, but together demonstrate that the
particle-pipeline machinery works correctly when used in its standard
form, while the restriction-via-bin-averaging shortcut does not.

\begin{figure}[!t]
\centering
\includegraphics[width=\textwidth]{zoom_ic_pk_full.pdf}
\caption{Power-spectrum diagnostics complementing
Fig.~\ref{fig:zoom-ic-cmp}. \textbf{Top row}: lo-res $P(k)$ of the
three restriction methods, measured on the coarse grid. A1 and A2
both develop a $W^2$-suppression deficit that grows toward the
coarse Nyquist; B (k-truncation) matches the input theory exactly.
\textbf{Bottom row}: standard particle-based IC validation on the
fine grid. 2LPT-displaced particles are deposited via 2D CIC; the
raw measurement is suppressed by the CIC window
$W_{\rm CIC}^2 = \prod_a\mathrm{sinc}^2(k_a\Delta x_f/2)$ and
deconvolution recovers the theory. Both rows use a fix-and-pair
white noise so per-mode cosmic-variance scatter is removed; what
remains is the systematic effect of each method or window. The two
rows live on different grids and answer different questions
(restriction error vs particle-pipeline window correction) but are
shown together because they share the underlying setup.}
\label{fig:zoom-ic-pk}
\end{figure}

\paragraph{What B still doesn't fix.}
Even approach B has a genuine boundary ``gap'' --- the high-$k$ modes
($|\mathbf k|>k_{\rm Ny,N}$) carried by the hi-res zoom interior simply
cannot be represented on the lo-res grid outside. This is a hard
Nyquist limit, not an algorithmic defect. Production zoom-IC codes
mitigate it in two complementary ways:
\begin{enumerate}
\item Surround the zoom region with a \emph{buffer zone} of additional
      hi-res cells. The gap amplitude decays with distance from the
      boundary on a scale of $\sim 1/k_{\rm Ny,N}$ (one coarse cell),
      so the science region --- which sits further inside the buffer
      --- is unaffected.
\item Use a multigrid (or hybrid real-space-multigrid plus FFT)
      Poisson solver that couples all grid levels in a single
      self-consistent solve, rather than independently FFT-Poisson
      per level. This avoids spurious mismatches in the gravitational
      potential at the boundary that would otherwise arise from
      treating each level's box as periodic in isolation.
\end{enumerate}
Both are remedies built on top of the basic observation that the
lo-res region must be constructed via k-truncation (B), not via
bin-averaging $\delta$ or $\phi^{(1)}$ (A1, A2).

\section{Summary}

\begin{itemize}
\item Block-averaging restriction is a smoothing step (convolution with
      a small averaging kernel) followed by $2{:}1$ subsampling. In
      Fourier space the smoothing multiplies by the cell window
      $W(\mathbf k)=\prod_a\cos(k_a\Delta x_f/2)$, and the subsampling
      folds fine-grid power from beyond the coarse Nyquist wavenumber
      back onto the coarse modes (Eq.~\ref{eq:restrict-fourier},
      Fig.~\ref{fig:density}).
\item The 1LPT displacement $\boldsymbol{\Psi}^{(1)}=i\mathbf k/k^2\,\hat\delta$
      is large-scale dominated and is therefore much more robust to
      restriction (Fig.~\ref{fig:psi}).
\item The deformation tensor $\Psi^{(1)}_{i,j}$ has kernel
      $-k_ik_j/k^2$ (bounded by unity) and shows a modest restriction
      error. The 2LPT source $S^{(2)}$, being quadratic in
      $\Psi^{(1)}_{i,j}$, inherits a noticeably larger error
      (Fig.~\ref{fig:deformation}). This is the practically
      interesting cost of restriction in an LPT pipeline.
\item Restricting $\phi^{(1)}$ and then differentiating produces a
      $\boldsymbol{\Psi}$ whose amplitude can be close to truth, but
      $-\nabla^2\phi^{(1)}_N\ne\delta_N$. The result is not the 1LPT
      displacement of the restricted density (Fig.~\ref{fig:phi}).
      When working across multiple grid resolutions one should therefore
      restrict the source ($\delta$, or the residual of the Poisson
      equation in iterative solvers) and re-solve on the coarser grid,
      rather than restrict the potential itself.
\end{itemize}

\appendix

\section{Convolution theorem in continuous and discrete settings}\label{app:convthm}

Several places in the main text use the fact that convolution in real
space becomes multiplication in Fourier space. This is the
\emph{convolution theorem}. The argument is essentially the same in
the continuous and discrete cases; we record both here.

\paragraph{Continuous case.}
Define the continuous Fourier transform with the convention
$\hat g(k) = \int g(x)\,e^{-ikx}\,dx$, and the convolution
\begin{equation}
  y(x) \;=\; (h*f)(x) \;\equiv\; \int h(x-x')\,f(x')\,dx'.
\end{equation}
Take the Fourier transform of $y$ and substitute the convolution
formula:
\begin{equation}
  \hat y(k) \;=\; \int y(x)\,e^{-ikx}\,dx
            \;=\; \iint h(x-x')\,f(x')\,e^{-ikx}\,dx\,dx'.
\end{equation}
Change variables $u = x-x'$ (so $x = u+x'$, $dx = du$), which makes
the exponent factor as $e^{-ikx} = e^{-iku}\,e^{-ikx'}$:
\begin{equation}
  \hat y(k) \;=\; \iint h(u)\,f(x')\,e^{-iku}\,e^{-ikx'}\,du\,dx'
            \;=\; \Bigl(\int h(u)\,e^{-iku}\,du\Bigr)
                  \Bigl(\int f(x')\,e^{-ikx'}\,dx'\Bigr).
\end{equation}
The two factors are just $\hat h(k)$ and $\hat f(k)$, so
\begin{equation}\label{eq:convthm-cont}
  \boxed{\;\hat y(k) \;=\; \hat h(k)\,\hat f(k).\;}
\end{equation}

\paragraph{Discrete case.}
Now the inputs are sequences $h_n, f_n$ on a grid of spacing
$\Delta x$. Define the continuous-$k$ Fourier transform of a sequence
by $\hat g(k) = \sum_n g_n\,e^{-ikn\Delta x}$ (the same formula used
for $H(k)$ in the main text), and the discrete convolution
\begin{equation}
  y_m \;=\; (h*f)_m \;\equiv\; \sum_n h_n\,f_{m-n}.
\end{equation}
Take the Fourier transform of $y$:
\begin{equation}
  \hat y(k) \;=\; \sum_m y_m\,e^{-ikm\Delta x}
            \;=\; \sum_{m,n} h_n\,f_{m-n}\,e^{-ikm\Delta x}.
\end{equation}
Substitute $m' = m-n$ (so $m = m'+n$, exponent
$e^{-ikm\Delta x} = e^{-ikn\Delta x}\,e^{-ikm'\Delta x}$):
\begin{equation}
  \hat y(k) \;=\; \sum_{n,m'} h_n\,f_{m'}\,e^{-ikn\Delta x}\,e^{-ikm'\Delta x}
            \;=\; \Bigl(\sum_n h_n\,e^{-ikn\Delta x}\Bigr)
                  \Bigl(\sum_{m'} f_{m'}\,e^{-ikm'\Delta x}\Bigr),
\end{equation}
which gives the same product law,
\begin{equation}\label{eq:convthm-disc}
  \boxed{\;\hat y(k) \;=\; \hat h(k)\,\hat f(k).\;}
\end{equation}
The substitution $m \to m'=m-n$ is the discrete analogue of the change
of variables $x \to u=x-x'$ in the continuous proof, and the rest of
the argument is identical: the double sum factors into a product of
single sums.

\paragraph{Convolution vs cross-correlation.}
The main text wrote the smoothing step as $y_m=\sum_n h_n\,f_{m+n}$
(a moving average), which is technically the discrete
\emph{cross-correlation} of $h$ with $f$, not the convolution. The two
differ by a sign-flip in $n$. For a kernel that is symmetric about its
support (the \emph{centered} version of our 2-sample kernel, with
weights $\tfrac12$ at $x=\pm\Delta x_f/2$, satisfies that property)
the two coincide, so there is no practical difference. The literal
sequence $h[0]=h[1]=\tfrac12$ used in the main text is not symmetric
about $n=0$, but its FT differs from the centered version's only by
the centroid-shift phase $e^{-ik\Delta x_f/2}$ discussed in
\S\ref{sec:fourier}, which drops out of any magnitude or
power-spectrum-level statement.

\section{Derivation of the alias-sum formula}\label{app:aliassum}

Equation~(\ref{eq:restrict-fourier}) is the central Fourier-space
identity for restriction. Here we derive it by composing the two
restriction steps in real space and translating each into Fourier
space, tracking the normalization prefactor explicitly.

\paragraph{Convention.}
We use the unnormalized continuous-$k$ Fourier transform of a discrete
sequence (consistent with App.~A): for a sequence $g_n$ on a grid of
spacing $\Delta x$,
\begin{equation}
  \hat g(k) = \sum_n g_n\,e^{-ikn\Delta x}.
\end{equation}
This is the convention that makes the convolution
theorem~(\ref{eq:convthm-disc}) read as a clean product, and it is the
convention against which the prefactor below is computed.

\paragraph{Step 1 --- smoothing.}
The first step is convolution of the fine-grid samples with the
averaging kernel $h$. By the convolution theorem
(\ref{eq:convthm-disc}), this becomes multiplication in Fourier space:
\begin{equation}\label{eq:app-smooth}
  \hat y(\mathbf k) \;=\; W(\mathbf k)\,\hat\delta_{2N}(\mathbf k),
  \qquad W(\mathbf k) \equiv \hat h(\mathbf k).
\end{equation}
$\hat y$ still lives on the fine grid --- same wavenumbers as
$\hat\delta_{2N}$.

\paragraph{Step 2 --- subsampling.}
For a sequence $g_n$ on the fine grid (spacing $\Delta x_f$), the
subsampled sequence $g'_I=g_{2I}$ on the coarse grid (spacing
$\Delta x_c=2\Delta x_f$) has the Fourier transform (this is the
standard ``decimation'' identity for Fourier transforms of discrete
sequences, derivable in one line from the even-mask trick
$\mathbb 1[n\,\mathrm{even}]=\tfrac12(1+e^{i\pi n})$):
\begin{equation}\label{eq:app-subsample-1d}
  \hat g'(k) \;=\; \frac{1}{2}\Bigl[\hat g(k) + \hat g(k + \pi/\Delta x_f)\Bigr]
  \qquad \text{(1D)}.
\end{equation}
Generalizing to $d$ dimensions with $2{:}1$ subsampling along each
axis, the prefactor becomes $1/2^d$ and the sum runs over $2^d$
shifts:
\begin{equation}\label{eq:app-subsample}
  \hat g'(\mathbf k) \;=\; \frac{1}{2^d}
  \sum_{\mathbf n\in\{0,1\}^d}
  \hat g\!\left(\mathbf k + \pi\mathbf n/\Delta x_f\right).
\end{equation}
Equivalently, in real space subsampling is multiplication by a Dirac
comb of spacing $\Delta x_c$, which Fourier-transforms to another
Dirac comb of spacing $2\pi/\Delta x_c=\pi/\Delta x_f$, and convolving
$\hat g$ with that comb (over one fine-mode period) gives the same
sum.

\paragraph{Combine.}
Applying~(\ref{eq:app-subsample}) to $g=y$ and substituting
$\hat y = W\hat\delta_{2N}$ from~(\ref{eq:app-smooth}):
\begin{equation}\label{eq:app-aliassum-rawconv}
  \hat\delta_N(\mathbf k_c) \;=\; \frac{1}{2^d}
  \sum_{\mathbf n\in\{0,1\}^d}
  W(\mathbf k_c + \pi\mathbf n/\Delta x_f)\,
  \hat\delta_{2N}(\mathbf k_c + \pi\mathbf n/\Delta x_f).
\end{equation}
The $\mathbf n=\mathbf 0$ term is the principal contribution
(multiplication by $W(\mathbf k_c)$ alone); the remaining $2^d-1$
terms (three in 2D, seven in 3D) are the alias contributions from
outer-wing fine modes that fold onto $\mathbf k_c$.

\paragraph{Reconciling with the main-text form.}
Equation~(\ref{eq:app-aliassum-rawconv}) is the alias-sum formula in
the unnormalized DTFT convention; it carries the prefactor $1/2^d$
(equivalently $1/M^d$ for $M{:}1$ subsampling). The main-text form
(\ref{eq:restrict-fourier}) drops this prefactor because the demo
scripts use the ``per-mode amplitude'' convention,
\begin{equation}
  \hat\delta_{2N}^{\rm pm}(\mathbf k) \equiv \hat\delta_{2N}(\mathbf k)/N_{2N}^d,
  \qquad
  \hat\delta_N^{\rm pm}(\mathbf k_c) \equiv \hat\delta_N(\mathbf k_c)/N_N^d,
\end{equation}
i.e., divide each grid's DFT by its total number of samples. Since
$N_{2N} = 2 N_N$, the relative normalization is $N_{2N}^d/N_N^d = 2^d$,
which exactly cancels the $1/2^d$ in
(\ref{eq:app-aliassum-rawconv}):
\begin{equation}
  \hat\delta_N^{\rm pm}(\mathbf k_c) \;=\;
  \sum_{\mathbf n\in\{0,1\}^d}
  W(\mathbf k_c + \pi\mathbf n/\Delta x_f)\,
  \hat\delta_{2N}^{\rm pm}(\mathbf k_c + \pi\mathbf n/\Delta x_f),
\end{equation}
which is (\ref{eq:restrict-fourier}). The choice of convention is a
matter of taste; the per-mode amplitude form is what NumPy's
\texttt{fftn} gives if you divide by the total grid size, and it is
the form most commonly used in cosmological power-spectrum estimation,
where one wants
$\langle|\hat\delta^{\rm pm}(\mathbf k)|^2\rangle = P(\mathbf k)/V$
to match the continuum normalization
$\langle\delta(\mathbf q_1)\delta(\mathbf q_2)\rangle = \int d^d k\,P(k)\,e^{i\mathbf k\cdot(\mathbf q_1-\mathbf q_2)}/(2\pi)^d$.

\section{Generating a Gaussian random field with prescribed $P(k)$}\label{app:grfgen}

The setup section uses the simplest of two equivalent ways to generate
a Gaussian random field with a target power spectrum $P(k)$ on a
periodic grid. Both produce the same statistical distribution; they
differ in code complexity and in how Hermitian symmetry is enforced.

\paragraph{Method A: white noise in real space, then FFT.}
\begin{enumerate}
\item Draw a unit-variance white-noise field
      $\mu(\mathbf q)\sim\mathcal N(0,1)$, independent and identically
      distributed (iid: same distribution at every site, statistically
      uncorrelated between sites) at every grid point.
\item Compute its Fourier transform $\hat\mu(\mathbf k)$.
\item Multiply mode-by-mode by $\sqrt{P(\mathbf k)}$ to colour the field:
      $\hat\delta(\mathbf k) = \sqrt{P(\mathbf k)}\,\hat\mu(\mathbf k)$.
\item Inverse-transform to real space; zero the
      $\mathbf k=\mathbf 0$ mode if you want exactly zero mean.
\end{enumerate}
Because $\mu$ is real, $\hat\mu(\mathbf k)$ is automatically
Hermitian-symmetric ($\hat\mu(-\mathbf k)=\hat\mu^*(\mathbf k)$) and
the inverse transform of $\hat\delta$ is therefore guaranteed real.
Special-case modes (DC and the Nyquist self-conjugate corners) are
real automatically.

\paragraph{Method B: white noise in Fourier space directly.}
\begin{enumerate}
\item Draw a complex Gaussian Fourier coefficient at each independent
      mode:
      \begin{equation}
        \hat\delta(\mathbf k) = \sqrt{P(\mathbf k)/2}\,(u + iv),
        \qquad u,v\sim\mathcal N(0,1)\;\text{iid}.
      \end{equation}
      With this normalization,
      $\mathrm{Re}\,\hat\delta$ and $\mathrm{Im}\,\hat\delta$ each have
      variance $P/2$, and $\langle|\hat\delta|^2\rangle = P$.
\item Equivalently in compact notation,
      $\hat\delta(\mathbf k) = \sqrt{P(\mathbf k)}\,z$ with
      $z\sim\mathcal{CN}(0,1)$, where the complex standard normal is
      $\mathcal{CN}(0,1)\equiv (u+iv)/\sqrt 2$ with $u,v\sim\mathcal N(0,1)$ iid.
\item Enforce Hermitian symmetry by hand:
      $\hat\delta(-\mathbf k) = \hat\delta^*(\mathbf k)$.
\item For self-conjugate modes (DC; Nyquist corners in even-sized
      grids), draw a real Gaussian rather than a complex one (no
      imaginary part).
\item Inverse-transform to real space.
\end{enumerate}
Method B requires more bookkeeping but lets you set the seed of
individual modes (useful for ``fixed-and-paired'' simulations or for
zoom ICs) and saves one FFT.

\paragraph{Equivalence on the grid.}
On an infinite (or sufficiently large) box, the two methods produce
statistically identical fields with
$\langle\hat\delta(\mathbf k)\hat\delta^*(\mathbf k')\rangle = P(\mathbf k)\,\delta_{\mathbf k,\mathbf k'}$.
On a \emph{finite} periodic box, however, they differ in their
treatment of long-wavelength modes. Method~A produces a field whose
realised spectrum is the box-window-convolved $P(\mathbf k)\otimes|\tilde W_C(\mathbf k)|^2$ ---
each grid mode inherits contributions from neighbouring (and even
sub-fundamental) continuous-$k$ modes through the sinc envelope of the
box window. Method~B assigns each grid mode the input $P$ value
exactly, with no sub-fundamental power at all (and enforces
$P(\mathbf k=\mathbf 0)=0$). The discrepancy matters most for small
boxes where a non-negligible fraction of the variance lies near or
below $k_{\rm fund}=2\pi/L$; \citet{Sirko2005} traced a $\sigma_8$
bias in $P$-sampled small-box ICs to exactly this effect. See
\texttt{notes/ic\_sampling.tex} and references therein
\citep{Pen1997, Sirko2005, Hahn2011} for a thorough discussion of the
$\xi$-sampled (Method A) versus $P$-sampled (Method B) distinction
and the convolution picture proposed by \citet{Pen1997}. The choice between them is therefore not just code
simplicity --- it has a real effect on large scales for small boxes.
All figures in this note use Method~A; for the boxes we work in, the
distinction is small at the wavenumbers where restriction effects are
visible (well above $k_{\rm fund}$).

\bibliography{cosmo_ic}

\end{document}
